% GENERATED FILE. Edit canonical lessons, then run npm run generate:books.
% canonical-source-hash: fnv1a64:90597e92b5c5e753
% canonical-lessons: ML-C202-prerana, ML-C202-nirbandham, ML-C202-gunam, ML-C202-dosam, ML-C202-nyayam

\chapter{Motivation, Insistence, Good quality, Drawback, Fairness}
\label{ch:ml-motivation-insistence-good-quality-drawback-fairness}
\hlchaptermodality{202}

\begin{chapteropening}
\textbf{By the end of this chapter:} \emph{I can say a motivation, insistence, a good quality, a drawback and fairness.}

Complete the last lesson of chapter 202: I can say a motivation, insistence, a good quality, a drawback and fairness.
\end{chapteropening}

\section[\texorpdfstring{\ml{പ്രേരണ} (prēraṇa)}{പ്രേരണ (prēraṇa)}]{\ml{പ്രേരണ} (prēraṇa) — a motivation, an urging}

\label{lesson:ML-C202-prerana}



\begin{quote}
Before the new one: say the Malayalam for unfortunately, then the Malayalam for somehow.
\end{quote}

\subsection*{You'll want to know: \ml{പ്രേരണ}}
\textbf{\ml{പ്രേരണ}} — \emph{prēraṇa} — \textquotedblleft{}a motivation, an urging\textquotedblright{}.

\begin{cousinweb}
\textbf{\ml{പ്രേരണ}} (\emph{prēraṇa}) $\leftarrow$ Sanskrit \emph{prēraṇa}, \textquotedblleft{}driving on, urging\textquotedblright{}, taken over whole as a tatsama (CDIAL 9000).
\end{cousinweb}

\begin{grammarlens}[title={What you've built}]
One more word for this chapter.
\end{grammarlens}

\subsection*{Your turn}
\begin{itemize}
\raggedright
  \item \emph{Say it:} \emph{prēraṇa}
  \item \emph{Say it:} \emph{prēraṇa}, once more
  \item \emph{Say it:} say \emph{prēraṇa}
  \item \emph{Recall:} say the Malayalam for unfortunately, then the Malayalam for somehow, then say \emph{prēraṇa} again
\end{itemize}

\begin{tcolorbox}[breakable,colback=teal!4,colframe=teal!35!black,title={Before you move on}]
What does \ml{പ്രേരണ} mean? (\textquotedblleft{}A motivation, an urging\textquotedblright{}.) Say it once more.
\end{tcolorbox}
\section[\texorpdfstring{\ml{നിർബന്ധം} (nirbandhaṁ)}{നിർബന്ധം (nirbandhaṁ)}]{\ml{നിർബന്ധം} (nirbandhaṁ) — insistence, compulsion}

\label{lesson:ML-C202-nirbandham}



\begin{quote}
Before the new one: say the Malayalam for somehow, then the Malayalam for a motivation.
\end{quote}

\subsection*{You'll want to know: \ml{നിർബന്ധം}}
\textbf{\ml{നിർബന്ധം}} — \emph{nirbandhaṁ} — \textquotedblleft{}insistence, compulsion\textquotedblright{}.

\begin{cousinweb}
\textbf{\ml{നിർബന്ധം}} (\emph{nirbandhaṁ}) $\leftarrow$ Sanskrit \emph{nirbandha}, \textquotedblleft{}perseverance, insistence\textquotedblright{}, taken over whole as a tatsama, with Malayalam's noun ending \emph{-ṁ} (CDIAL 7355). It is built on the root \emph{bandh}, \textquotedblleft{}to bind\textquotedblright{}.
\end{cousinweb}

\begin{grammarlens}[title={What you've built}]
One more word for this chapter.
\end{grammarlens}

\subsection*{Your turn}
\begin{itemize}
\raggedright
  \item \emph{Say it:} \emph{nirbandhaṁ}
  \item \emph{Say it:} \emph{nirbandhaṁ}, once more
  \item \emph{Say it:} say \emph{nirbandhaṁ}
  \item \emph{Recall:} say the Malayalam for somehow, then the Malayalam for a motivation, then say \emph{nirbandhaṁ} again
\end{itemize}

\begin{tcolorbox}[breakable,colback=teal!4,colframe=teal!35!black,title={Before you move on}]
What does \ml{നിർബന്ധം} mean? (\textquotedblleft{}Insistence, compulsion\textquotedblright{}.) Say it once more.
\end{tcolorbox}
\section[\texorpdfstring{\ml{ഗുണം} (guṇaṁ)}{ഗുണം (guṇaṁ)}]{\ml{ഗുണം} (guṇaṁ) — a good quality, a virtue}

\label{lesson:ML-C202-gunam}



\begin{quote}
Before the new one: say the Malayalam for a motivation, then the Malayalam for insistence.
\end{quote}

\subsection*{You'll want to know: \ml{ഗുണം}}
\textbf{\ml{ഗുണം}} — \emph{guṇaṁ} — \textquotedblleft{}a good quality, a virtue\textquotedblright{}.

\begin{cousinweb}
\textbf{\ml{ഗുണം}} (\emph{guṇaṁ}) $\leftarrow$ Sanskrit \emph{guṇa}, \textquotedblleft{}strand; quality, good quality\textquotedblright{}, taken over whole as a tatsama, with Malayalam's noun ending \emph{-ṁ} (CDIAL 4190). Its first sense is a strand of a cord.
\end{cousinweb}

\begin{grammarlens}[title={What you've built}]
One more word for this chapter.
\end{grammarlens}

\subsection*{Your turn}
\begin{itemize}
\raggedright
  \item \emph{Say it:} \emph{guṇaṁ}
  \item \emph{Say it:} \emph{guṇaṁ}, once more
  \item \emph{Say it:} say \emph{guṇaṁ}
  \item \emph{Recall:} say the Malayalam for a motivation, then the Malayalam for insistence, then say \emph{guṇaṁ} again
\end{itemize}

\begin{tcolorbox}[breakable,colback=teal!4,colframe=teal!35!black,title={Before you move on}]
What does \ml{ഗുണം} mean? (\textquotedblleft{}A good quality, a virtue\textquotedblright{}.) Say it once more.
\end{tcolorbox}
\section[\texorpdfstring{\ml{ദോഷം} (dōṣaṁ)}{ദോഷം (dōṣaṁ)}]{\ml{ദോഷം} (dōṣaṁ) — a drawback, harm}

\label{lesson:ML-C202-dosam}



\begin{quote}
Before the new one: say the Malayalam for insistence, then the Malayalam for a good quality.
\end{quote}

\subsection*{You'll want to know: \ml{ദോഷം}}
\textbf{\ml{ദോഷം}} — \emph{dōṣaṁ} — \textquotedblleft{}a drawback, harm\textquotedblright{}.

\begin{cousinweb}
\textbf{\ml{ദോഷം}} (\emph{dōṣaṁ}) $\leftarrow$ Sanskrit \emph{dōṣa}, \textquotedblleft{}fault\textquotedblright{}, taken over whole as a tatsama, with Malayalam's noun ending \emph{-ṁ} (CDIAL 6587).
\end{cousinweb}

\begin{grammarlens}[title={What you've built}]
One more word for this chapter.
\end{grammarlens}

\subsection*{Your turn}
\begin{itemize}
\raggedright
  \item \emph{Say it:} \emph{dōṣaṁ}
  \item \emph{Say it:} \emph{dōṣaṁ}, once more
  \item \emph{Say it:} say \emph{dōṣaṁ}
  \item \emph{Recall:} say the Malayalam for insistence, then the Malayalam for a good quality, then say \emph{dōṣaṁ} again
\end{itemize}

\begin{tcolorbox}[breakable,colback=teal!4,colframe=teal!35!black,title={Before you move on}]
What does \ml{ദോഷം} mean? (\textquotedblleft{}A drawback, harm\textquotedblright{}.) Say it once more.
\end{tcolorbox}
\section[\texorpdfstring{\ml{ന്യായം} (nyāyaṁ)}{ന്യായം (nyāyaṁ)}]{\ml{ന്യായം} (nyāyaṁ) — fairness, a justification}

\label{lesson:ML-C202-nyayam}



\begin{quote}
Before the new one: say the Malayalam for a good quality, then the Malayalam for a drawback.
\end{quote}

\subsection*{You'll want to know: \ml{ന്യായം}}
\textbf{\ml{ന്യായം}} — \emph{nyāyaṁ} — \textquotedblleft{}fairness, a justification\textquotedblright{}.

\begin{cousinweb}
\textbf{\ml{ന്യായം}} (\emph{nyāyaṁ}) $\leftarrow$ Sanskrit \emph{nyāya}, \textquotedblleft{}rule, method\textquotedblright{}, taken over whole as a tatsama, with Malayalam's noun ending \emph{-ṁ} (CDIAL 7614).
\end{cousinweb}

\begin{grammarlens}[title={What you've built}]
One more word for this chapter.
\end{grammarlens}

\subsection*{Your turn}
\begin{itemize}
\raggedright
  \item \emph{Say it:} \emph{nyāyaṁ}
  \item \emph{Say it:} \emph{nyāyaṁ}, once more
  \item \emph{Say it:} say \emph{nyāyaṁ}
  \item \emph{Recall:} say the Malayalam for a good quality, then the Malayalam for a drawback, then say \emph{nyāyaṁ} again
\end{itemize}

\begin{tcolorbox}[breakable,colback=teal!4,colframe=teal!35!black,title={Before you move on}]
What does \ml{ന്യായം} mean? (\textquotedblleft{}Fairness, a justification\textquotedblright{}.) Say it once more.
\end{tcolorbox}
